
        \documentclass[spanish, 11pt]{exam}

        %These tell TeX which packages to use.
        \usepackage{array,epsfig}
        \usepackage{amsmath, textcomp}
        \usepackage{amsfonts}
        \usepackage{amssymb}
        \usepackage{amsxtra}
        \usepackage{amsthm}
        \usepackage{mathrsfs}
        \usepackage{color}
        \usepackage{multicol, xparse}
        \usepackage{verbatim}


        \usepackage[utf8]{inputenc}
        \usepackage[spanish]{babel}
        \usepackage{eurosym}

        \usepackage{graphicx}
        \graphicspath{{../img/}}
        \usepackage{pgf}



        \printanswers
        \nopointsinmargin
        \pointformat{}

        %Pagination stuff.
        %\setlength{\topmargin}{-.3 in}
        %\setlength{\oddsidemargin}{0in}
        %\setlength{\evensidemargin}{0in}
        %\setlength{\textheight}{9.in}
        %\setlength{\textwidth}{6.5in}
        %\pagestyle{empty}

        \let\multicolmulticols\multicols
        \let\endmulticolmulticols\endmulticols
        \RenewDocumentEnvironment{multicols}{mO{}}
         {%
          \ifnum#1=1
            #2%
          \else % More than 1 column
            \multicolmulticols{#1}[#2]
          \fi
         }
         {%
          \ifnum#1=1
          \else % More than 1 column
            \endmulticolmulticols
          \fi
         }
        \renewcommand{\solutiontitle}{\noindent\textbf{Sol:}\enspace}

        \newcommand{\samedir}{\mathbin{\!/\mkern-5mu/\!}}

        \newcommand{\class}{1º Bachillerato}
        \newcommand{\examdate}{\today}

        \newcommand{\tipo}{A}


        \newcommand{\timelimit}{50 minutos}



        \pagestyle{head}
        \firstpageheader{\includegraphics[width=0.2\columnwidth]{header_left}}{\textbf{Departamento de Matemáticas\linebreak \class}\linebreak \examnum}{\includegraphics[width=0.1\columnwidth]{header_right}}
        \runningheader{\class}{\examnum}{Página \thepage\ of \numpages}
        \runningheadrule

        \newcommand{\examnum}{44 - Variables aleatorias}
        \begin{document}
        \begin{questions}
        \question p105e02-0 - Sea X una variable aleatoria discreta cuya función de probabilidad es ($x_i:p_i$):\\{0: 0.1, 1: 0.2, 2: 0.1, 3: 0.4, 4: 0.1, 5: 0.1} 
    
        \begin{multicols}{1}
        \begin{parts} \part[1] Calcula sus parámetros  \begin{solution}   \\La media es: $2.5$ \\La varianza: $2.05$   \end{solution} \part[1] Calcula $P\left(X < 4.5\right)$: , $P\left( X \geqslant 3 \right)$, $P\left(3\leqslant X < 4.5 \right)$   \begin{solution}    $\left [ 0.9, \quad 0.6, \quad 0.5\right ]$   \end{solution}
        \end{parts}
        \end{multicols}
        \question p105e08-0 - En una distribución binomial B(9 , 0.2) calcula:\\ 
    
        \begin{multicols}{1}
        \begin{parts} \part[1] $P\left(X > 3\right)$  \begin{solution}    $0.085641728$   \end{solution} \part[1] $P\left(X \geq 7\right)$  \begin{solution}    $0.000313856$   \end{solution} \part[1] $P\left(X > 0\right)$  \begin{solution}    $0.865782272$   \end{solution} \part[1] $P\left(X \leq 9\right)$  \begin{solution}    $1.0$   \end{solution}
        \end{parts}
        \end{multicols}
        \question p105e09-0 - La última novela de cierto afamado autor ha tenido un importante éxito, hasta el punto de que el
80\% de los lectores ya la han leído. Un grupo de 4 amigos son aficionados a la lectura\\ 
    
        \begin{multicols}{1}
        \begin{parts} \part[1] Describe la variable que indica el número de individuos del grupo que han leído la novela  \begin{solution}   $\left \{ 0 : 0.0016, \quad 1 : 0.0256, \quad 2 : 0.1536, \quad 3 : 0.4096, \quad 4 : 0.4096\right \}$   \end{solution} \part[1] ¿Cuál es la probabilidad de que en el grupo hayan leído la obra 2 personas? ¿Y al menos 2?  \begin{solution}   $0.1536$ y $0.9728$   \end{solution}
        \end{parts}
        \end{multicols}
        
    \end{questions}
    \end{document}
    